% Перечень сокращений и обозначений
\documentclass[../main.tex]{subfiles}

\begin{document}

\onlyinsubfile{
    \renewcommand{\contentsname}{ОГЛАВЛЕНИЕ}
    \setcounter{tocdepth}{3}
    \tableofcontents
}

\chapter*{\MakeUppercase{Перечень сокращений и обозначений}}
\addcontentsline{toc}{chapter}{Перечень сокращений и обозначений}

\begin{longtable}{p{0.2\textwidth}p{0.7\textwidth}}

AAE & Adversarial Autoencoder (состязательный автоэнкодер) \\
AUC & Area Under the Curve (площадь под кривой) \\
CL & Curriculum Learning (обучение по учебному плану) \\
CPU & Central Processing Unit (центральный процессор) \\
DAP & Difference between Augmented and Posterior \\
DFT & Density Functional Theory (теория функционала плотности) \\
ELBO & Evidence Lower Bound (нижняя граница правдоподобия) \\
ELC & Explicit Latent Conditioning (явное обусловливание латентного пространства) \\
FCD & Fréchet ChemNet Distance (расстояние Фреше в пространстве ChemNet) \\
GCN & Graph Convolutional Networks (графовые сверточные сети) \\
GNN & Graph Neural Networks (графовые нейронные сети) \\
GPU & Graphics Processing Unit (графический процессор) \\
GRU & Gated Recurrent Unit (управляемый рекуррентный блок) \\
InChI & International Chemical Identifier (международный химический идентификатор) \\
JT-VAE & Junction Tree Variational Autoencoder \\
LLM & Large Language Model (большая языковая модель) \\
LogP & Коэффициент распределения октанол-вода \\
LSTM & Long Short-Term Memory (долгая краткосрочная память) \\
MD-VAE & Multi-Decoder Variational Autoencoder \\
MOSES & Molecular Sets \\
NLP & Natural Language Processing (обработка естественного языка) \\
OOD & Out-of-Distribution (вне распределения) \\
pLogP & Penalized LogP (штрафованный LogP) \\
QSAR & Quantitative Structure-Activity Relationship (количественное соотношение структура-активность) \\
RL & Reinforcement Learning (обучение с подкреплением) \\
RNN & Recurrent Neural Network (рекуррентная нейронная сеть) \\
SA score & Synthetic Accessibility score (оценка синтетической доступности) \\
SMILES & Simplified Molecular-Input Line-Entry System (упрощенная система представления молекул) \\
TPSA & Topological Polar Surface Area (топологическая полярная площадь поверхности) \\
VAE & Variational Autoencoder (вариационный автоэнкодер) \\
АИ & Искусственный интеллект \\
ВИЧ & Вирус иммунодефицита человека \\
$\Delta G$ & Энергия связывания Гиббса \\
$EC_{50}$ & Концентрация полумаксимального эффекта \\
$R^2$ & Коэффициент детерминации \\

\end{longtable}


\end{document}
